% =====================================================================================================================
\section{Using \tool{}}
\label{sec:using}

\tool{} ships as a Claude Code plugin with a command-line tool inside it. This section describes what a user installs
and what happens in a session, with the choices our measurements support.

\paragraph{Setup.}
The plugin is installed once and stays silent until a project opts in. A project opts in by having a file
\code{paper/sources.py}, which \code{paperanchor init} creates:
\begin{lstlisting}[language=bash]
$ claude plugin marketplace add thanh-huy-nguyen/paperanchor
$ claude plugin install paperanchor@paperanchor
$ cd my-project && paperanchor init     # creates paper/sources.py and snapshots the repository
\end{lstlisting}
The paper's preamble gains two lines that make the tags invisible in the typeset output. Nothing else about the
paper changes.
\begin{lstlisting}
\providecommand{\src}[1]{}       % tag after a number: prints nothing
\providecommand{\val}[2][1]{}    % placeholder filled by `paperanchor fill`
\end{lstlisting}

\paragraph{What happens in a session.}
At session start, the plugin snapshots the repository and tells the agent the convention in a few lines, so the user
does not have to put it in a prompt. A bundled skill teaches the agent how to write recipes. From then on every edit
to \code{paper/*.tex} is checked (Table~\ref{tab:events}). The agent writes the recipe first and the number second,
and a number that does not follow comes back with the value the data gives.

\begin{table}[t]
  \centering
  \caption{What \tool{} does at each point of a Claude Code session, by mode. Feedback is the default.}
  \label{tab:events}
  \small
  \setlength{\tabcolsep}{5pt}
  \begin{tabular}{@{}p{0.27\linewidth}p{0.32\linewidth}p{0.33\linewidth}@{}}
    \toprule
    \rowcolor{headgray}
    Event & Feedback mode & Gate mode \\
    \midrule
    Session start & \multicolumn{2}{l}{Snapshot the repository; tell the agent the convention} \\
    Edit or write to the paper & Keep the edit; report every failing number & Deny the edit if a number it adds or
    changes fails \\
    Shell command that changes the paper & Keep the change; report failures & Revert the change \\
    Agent stops & Nothing & Keep the session going while a tagged number fails, up to 3 times \\
    Commit (git hook) & \multicolumn{2}{l}{Recheck every tagged number in staged papers; refuse the commit on a failure} \\
    \bottomrule
  \end{tabular}
\end{table}

\paragraph{Which mode to use.}
Feedback is the default because it reached the gate's outcome with fewer than half the turns
(Table~\ref{tab:placement}). The gate is for sessions that run unattended and must not end with a failing number.
Generate mode suits result-heavy sections: the agent writes \valmacro{} and \code{paperanchor fill} types the number,
so a rerun of the experiments updates the paper without another writing session. A log mode records every decision
and shows the agent nothing, which is how we ran the arms without a check.

\paragraph{What it costs.}
The check makes no model call. A full recheck of a paper reruns every recipe; on the eight papers of the placement
study, which carry between 21 and 194 tagged numbers (Appendix Table~\ref{tab:repos}), it took \checktime{}. The cost a user feels is in the writing: the
agent spends a few lines of Python on each number, and in our study it wrote fewer new numbers when it had to justify
each one (Table~\ref{tab:placement}).

\paragraph{Outside Claude Code.}
The command-line tool does not need an agent. \code{paperanchor check} verifies every tagged number of a paper and
exits non-zero on a failure, so it can run in continuous integration, and \code{paperanchor git-hook} installs it as a
pre-commit hook. A co-author who edits a number by hand is held to the same recipes as the agent.

\paragraph{What to watch for.}
Three engineering failures surfaced in our early runs, and each made a broken hook look like a passing one: a hook
that exceeded its timeout did nothing, an agent that changed directory made the hook look for the paper in the wrong
place, and a hook that covered only the edit tools missed an agent that edited the paper through the shell. \tool{}
now logs every decision, resolves the project root from the edited file, and covers shell commands. A check that can
fail silently should be assumed to have done so until its log says otherwise.
